\section{The Position of \VED{}}

\label{sec:ved}

This volume introduced the structure of the \DH{}
and reinterpreted divergence, renormalization,
and the limits of modern physical theory.

Within this context,
the position of \VED{} (\textit{Vortical Enclosure Dynamics})
can be clarified.

\medskip
\begin{quote}
\VED{} is not a unified theory.
\end{quote}
\medskip

It does not aim to integrate all interactions
into a single framework.

\medskip
\begin{quote}
\VED{} does not describe the \ungenerated{} domain.
\end{quote}
\medskip

The \ungenerated{} is in principle indescribable.
It is not an object of theory,
but appears only as a boundary.

\medskip
\begin{quote}
\VED{} may be understood as a theory of the boundary of generation.
\end{quote}
\medskip

It describes the structure at the \DH{}
and the behavior of theory in its vicinity.

\medskip
\begin{quote}
\VED{} does not close theory.
\end{quote}
\medskip

The \DH{} cannot be reached.
Therefore,
no theory can fully describe it.

\medskip
\begin{quote}
This limit is not a deficiency,\\
but the condition under which theory can remain consistent.
\end{quote}
\medskip

\VED{} does not replace existing theories.

Quantum field theory, relativity, and string theory
remain valid within their respective domains.

\medskip
\begin{quote}
\VED{} shows the structure within which these theories are positioned.
\end{quote}
\medskip

\medskip
\begin{quote}
Theory does not end here.\\
It simply cannot proceed further in the same form.
\end{quote}
\medskip

The \DH{} lies outside all theories.
Yet all theories approach it.

\medskip
\begin{quote}
Physics operates within the horizon.\\
Of what lies beyond, nothing can be said.
\end{quote}

\begin{quote}
What lies beyond the boundary remains unsaid.
\end{quote}
\medskip

This volume has not provided a complete description of the world.

\medskip
\begin{quote}
What has been reached is not an answer,\\
but a boundary.
\end{quote}
